% Mathematical source: ../1-practice_activities/G11_T2_L1_C2_practice.tex
% User override: one worked, one guided, four independent problems (1--6).
\makeatletter
\def\input@path{{../../../preamble/}}
\makeatother
\input{preamble.tex}

% Highlight infrastructure is used only by the appended solutions.
\input{graphs/highlights.tex}
% \GTenHighlightsfalse % Uncomment to disable instructional highlights.
\input{graphs/specific_highlights.tex}
\newcommand{\CtwoSignificance}[1]{\CtwoConfidence{#1}}
\definecolor{CtwoErrorColor}{HTML}{28AFA2}
\newcommand{\CtwoError}[1]{\HighlightStroke{CtwoErrorColor}{#1}}
\newcommand{\CtwoCI}[1]{\CtwoEndpoint{#1}}
\newcommand{\CtwoCoverCount}{\CtwoCoverage{N_{\rm cover}}}
\newcommand{\CtwoTotalCount}{\CtwoCoverage{N}}

\newcommand{\SolutionInformationTable}[1]{%
    \begingroup
    \small
    \setlength{\tabcolsep}{4pt}
    \renewcommand{\arraystretch}{1.16}
    \noindent
    \begin{tabularx}{\linewidth}{@{}>{\raggedright\arraybackslash}p{0.08\linewidth}>{\raggedright\arraybackslash}p{0.18\linewidth}>{\raggedright\arraybackslash}p{0.18\linewidth}>{\raggedright\arraybackslash}X@{}}
        \toprule
        \textbf{Symbol} & \textbf{Concept} & \textbf{Type} & \textbf{Meaning in context} \\
        \midrule
        #1
        \bottomrule
    \end{tabularx}\par
    \vspace{2mm}
    \endgroup
}

\newcommand{\SolutionTaskTwoTable}[1]{%
    \par\begingroup
    \small
    \setlength{\tabcolsep}{4pt}
    \renewcommand{\arraystretch}{1.18}
    \noindent
    \begin{tabularx}{\linewidth}{@{}>{\raggedright\arraybackslash}p{0.13\linewidth}>{\raggedright\arraybackslash}p{0.16\linewidth}>{\raggedright\arraybackslash}p{0.28\linewidth}>{\raggedright\arraybackslash}p{0.19\linewidth}>{\raggedright\arraybackslash}X@{}}
        \toprule
        \textbf{Symbol} & \textbf{Concept} & \textbf{Calculation} & \textbf{Result} & \textbf{Meaning in context} \\
        \midrule
        #1
        \bottomrule
    \end{tabularx}\par\vspace{2mm}
    \endgroup
}

\usepackage{needspace}
\usepackage{varwidth}
\geometry{letterpaper,margin=0.25in}
\pagestyle{empty}
\linespread{1.04}
\setlength{\parskip}{4pt}
\raggedbottom
\newcommand{\stageheading}[1]{\par\addvspace{0.5\baselineskip}{\centering\large\bfseries #1\par}\nopagebreak[4]}

\newcommand{\problemtitle}[2]{\par\Needspace{4\baselineskip}\addvspace{.4\baselineskip}\noindent\textbf{#1. #2}\par\nopagebreak[4]}
\newcommand{\workspace}[1]{\par\vspace{#1}}
\newcommand{\task}[2]{\par\Needspace{3\baselineskip}\textbf{Task #1 -- #2}}
\newcommand{\InformationTask}{\textbf{Task 1 -- Identify and organize the information; calculate SE.} In independent practice, use four columns: \textbf{Symbol}, \textbf{Concept}, \textbf{Type}, and \textbf{Meaning in context}. The examples use two side-by-side three-column tables, with the categories as section titles instead of a Type column. Arrange three contiguous groups in this order: \textbf{Population parameters} ($\mu$, $\mu_{\mathrm{known}}$, $\sigma$); \textbf{Sample statistics} ($\bar{x}$, $n$, $\mathrm{SE}$); \textbf{Procedure Settings} ($C$, $\alpha$, $\alpha/2$, $z^*$).}
\newcommand{\IntervalTask}{\textbf{Task 2 -- Calculate MoE and state the interval.} In independent practice, use four columns: \textbf{Symbol}, \textbf{Concept}, \textbf{Calculation}, and \textbf{Result and meaning in context}. In both examples, combine the last two columns as \textbf{Calculation, Result and meaning in context}. Show formulas and substitutions, then results with units and brief context. Use rows for the point estimate, margin of error, confidence interval, and full width. In independent practice, also use separate endpoint rows. In the examples, include endpoint calculations in the confidence-interval row. Group by confidence level; give the shared point estimate once. Keep all calculated endpoint digits. Calculate the full width as $W=U-L=2\mathrm{MoE}$.}
\newcommand{\CoverageTask}{\textbf{Task 3 -- Compare observed interval coverage with nominal confidence.} For each level separately, count subsequent intervals satisfying $L_i\leq\mu_{\mathrm{known}}\leq U_i$; divide by the total number of intervals and multiply by $100\%$. Compare with the corresponding nominal confidence percentage used in the initial study and calculate the absolute difference in percentage points. Interpret discrepancies as finite-sample variation under the stated model: coverage need not equal nominal confidence, especially with few studies, and need not approach it monotonically as studies are added.}

\newcommand{\InformationTable}[1]{%
    \par
    \begingroup
    \small
    \setlength{\tabcolsep}{4pt}
    \renewcommand{\arraystretch}{1.16}
    \noindent
    \begin{tabularx}{\linewidth}{@{}>{\raggedright\arraybackslash}p{0.10\linewidth}>{\raggedright\arraybackslash}p{0.24\linewidth}>{\raggedright\arraybackslash}p{0.24\linewidth}>{\raggedright\arraybackslash}X@{}}
        \toprule
        \textbf{Symbol} & \textbf{Concept} & \textbf{Type} & \textbf{In context...} \\
        \midrule
        #1
        \bottomrule
    \end{tabularx}\par
    \vspace{2mm}
    \endgroup
}

% Each confidence-level block stays together; repeat headings in every block.
\newcommand{\EstimationTable}[1]{%
    \par
    \begingroup
    \setlength{\tabcolsep}{4pt}
    \renewcommand{\arraystretch}{1.18}
    \noindent
    \begin{tabularx}{\linewidth}{@{}>{\raggedright\arraybackslash}p{0.095\linewidth}>{\raggedright\arraybackslash}p{0.17\linewidth}>{\raggedright\arraybackslash}p{0.27\linewidth}>{\raggedright\arraybackslash}X@{}}
        \toprule
        \textbf{Symbol} & \textbf{Concept} & \textbf{Calculation} & \textbf{Result and meaning in context} \\
        \midrule
        #1
        \bottomrule
    \end{tabularx}\par
    \vspace{2mm}
    \endgroup
}


% Dedicated example tables; independent-practice macros remain unchanged.
\newcommand{\ExampleInformationTable}[1]{%
    \par
    \begingroup
    \small
    \setlength{\tabcolsep}{4pt}
    \renewcommand{\arraystretch}{1.16}
    \noindent\begin{tabularx}{\linewidth}{|>{\raggedright\arraybackslash}p{.14\linewidth}|>{\raggedright\arraybackslash}p{.32\linewidth}|>{\raggedright\arraybackslash}X|}
        \hline
        \textbf{Symbol} & \textbf{Concept} & \textbf{In context...} \\ \hline
        #1
    \end{tabularx}\par
    \endgroup
}
\newcommand{\GuidedEstimationTable}[2][.1\linewidth]{%
    \par
    \begingroup
    \small
    \setlength{\tabcolsep}{4pt}
    \renewcommand{\arraystretch}{1.16}
    \noindent\begin{tabularx}{\linewidth}{|>{\raggedright\arraybackslash}p{#1}|>{\raggedright\arraybackslash}p{.2\linewidth}|>{\raggedright\arraybackslash}X|}
        \hline
        \textbf{Symbol} & \textbf{Concept} & \textbf{Calculation, Result and meaning in context} \\ \hline
        #2
    \end{tabularx}\par
    \endgroup
}
\newcommand{\GuidedConfidenceTable}[1]{%
    \GuidedEstimationTable[.14\linewidth]{
        \multicolumn{3}{|l|}{\textbf{#1 confidence}}\\ \hline
        $\mathrm{MoE}$ & Margin of error & \\[10pt] \hline
        $\mathrm{CI}$ & Confidence interval & \\[10pt] \hline
        $W$ & Full width & \\[10pt] \hline
    }
}

\newcommand{\BlankInformationTable}{
\InformationTable{
$\mu$&Population mean&Population parameters&\\[23pt]
$\mu_{\mathrm{known}}$&Later established population mean&&\\[18pt]
$\sigma$&Population standard deviation&&\\[18pt]\midrule
$\bar{x}$&Sample mean&Sample statistics&\\[18pt]
$n$&Sample size&&\\[18pt]
$\mathrm{SE}$&Standard error&&\\[27pt]\midrule
$C$&Confidence level&Procedure Settings&\\[18pt]
$\alpha$&Significance level&&\\[18pt]
$\alpha/2$&Tail area&&\\[18pt]
$z^*$&Critical value&&\\[18pt]}}
\begin{document}
{\centering{\Large\bfseries C2 -- Point and interval estimation\par}
Global Economics \textbar{} Grade 11 \textbar{} Term 2 \textbar{} Lesson 1\par}
\textbf{C2:} Compare point and interval estimation and interpret confidence through long-run coverage, distinguishing this from the probability that a particular calculated interval contains the parameter.

\stageheading{Equations and basic procedure}
Assume independent random observations, known $\sigma$, and normally distributed $\bar{x}$ (exact for normal populations). Each problem's studies sample the same fixed population for the same target period, using the same $n$, known $\sigma$, and confidence-interval procedure at each stated level. Initially estimate unknown $\mu$ with $\bar{x}$. A later comprehensive population investigation establishes $\mu_{\mathrm{known}}$, the value of that same parameter, treated as known for this activity; repeated sample estimates alone do not establish it. No sample SD $s$ is needed.

Each subsequent study constructs its own interval $[L_i,U_i]$ around its own sample mean $\bar{x}_i$. It covers the population mean when $L_i\leq\mu_{\mathrm{known}}\leq U_i$. The nominal confidence level describes the procedure's long-run interval coverage, including the procedure used for the initial interval. It is neither the percentage of subsequent sample means nor individual observations inside the initial interval, nor a probability assigned to the fixed mean after calculation.

\begingroup
\small
\setlength{\tabcolsep}{3pt}
\setlength{\extrarowheight}{0pt}
\setlength{\parskip}{0pt}
\renewcommand{\arraystretch}{1}
\noindent\begin{minipage}[t]{.49\linewidth}
\vspace{0pt}
\begin{tabularx}{\linewidth}{@{}>{\raggedright\arraybackslash}p{.21\linewidth}>{\raggedright\arraybackslash}p{.34\linewidth}>{\raggedright\arraybackslash}X@{}}
\toprule
\textbf{Symbol} & \textbf{Concept} & \textbf{Formula} \\
\midrule
$C$ & Nominal confidence level & \\
$\mu$ & Population mean & Initially unknown \\
$\mu_{\mathrm{known}}$ & Later established population mean & Same parameter $\mu$ \\
$z^*$ & Critical value & \\
$\mathrm{SE}$ & Standard error & $\mathrm{SE}=\frac{\sigma}{\sqrt n}$ \\
$\mathrm{MoE}$ & Margin of error & $\mathrm{MoE}=z^*\mathrm{SE}$ \\
$L$ & Lower endpoint & $L=\bar{x}-\mathrm{MoE}$ \\
$U$ & Upper endpoint & $U=\bar{x}+\mathrm{MoE}$ \\
\bottomrule
\end{tabularx}\par
\end{minipage}\hfill
\begin{minipage}[t]{.49\linewidth}
\vspace{0pt}
\begin{tabularx}{\linewidth}{@{}>{\raggedright\arraybackslash}p{.21\linewidth}>{\raggedright\arraybackslash}p{.34\linewidth}>{\raggedright\arraybackslash}X@{}}
\toprule
\textbf{Symbol} & \textbf{Concept} & \textbf{Formula} \\
\midrule
$\mathrm{CI}$ & Confidence interval & $\mathrm{CI}=[L,U]$ \\
$W$ & Full width & $W=U-L=2\mathrm{MoE}$ \\
$\alpha$ & Significance level & $\alpha=1-C$ \\
$\alpha/2$ & Each tail's area & $\frac{\alpha}{2}=\frac{1-C}{2}$ \\
$N_{\rm cover}(C)$ & Covering count & \\
$N(C)$ & Total interval count & \\
$\mathrm{Cov}(C)$ & Observed coverage & $\frac{N_{\rm cover}(C)}{N(C)}\times100\%$ \\
\bottomrule
\end{tabularx}\par
\end{minipage}\par
\endgroup
\noindent\begin{minipage}[t]{0.68\linewidth}
\vspace{0pt}
\[
\begin{aligned}
\mathrm{Cov}(C)&=\frac{N_{\rm cover}(C)}{N(C)}\times100\%,\\
E_{\rm abs}(C)&=|\mathrm{Cov}(C)-C|.
\end{aligned}
\]
Use coverage and nominal confidence as percentages here; $E_{\rm abs}$ is their absolute difference in percentage points. For tail areas, use $C$ as a proportion.
\end{minipage}\hfill
\begin{minipage}[t]{0.30\linewidth}
\vspace{0pt}
\centering
\small
\textbf{Two-sided confidence levels and standard normal critical values}\par
\begin{tabularx}{\linewidth}{@{}>{\raggedright\arraybackslash}X>{\raggedleft\arraybackslash}X@{}}\toprule
Confidence level $C$ & Critical value $z^*$\\\midrule
90\% & 1.645\\
95\% & 1.960\\
99\% & 2.576\\
99.5\% & 2.807\\
99.9\% & 3.291\\
\bottomrule
\end{tabularx}
\end{minipage}\par



\vspace{-4mm}

\stageheading{General assessment tasks}

\InformationTask

\IntervalTask

\CoverageTask

Legacy labels in the preserved solutions refer to $\mu_{\mathrm{known}}$, the later established mean of the same population.

\clearpage
\stageheading{Worked example}
\vspace{-7mm}

\problemtitle{1}{Investigate the cycling interval}
\noindent\begin{minipage}[t]{0.68\linewidth}
\vspace{0pt}
An initial study estimated the then unknown average duration of all cycling trips in a specified month. A random sample of 49 trips averaged 40 minutes; the known population standard deviation was 14 minutes. The report used 95\% confidence.

Five subsequent independent studies sampled the same month's population, each using 49 trips, the same standard deviation and 95\% procedure. A later comprehensive population investigation established an average of 41 minutes, enabling retrospective interval assessment.
\end{minipage}\hfill
\begin{minipage}[t]{0.30\linewidth}
\vspace{0pt}
\centering
\begin{tabular}{@{}rrrr@{}}\toprule
Study&$\bar{x}$ (min)& $L$(min)&$U$ (min)\\\midrule
1&40&36.08&43.92\\
2&42&38.08&45.92\\
3&45&41.08&48.92\\
4&38&34.08&41.92\\
5&41&37.08&44.92\\\bottomrule
\end{tabular}
\end{minipage}\par

\vspace{-3mm}
\subsubsection*{Task 1 solution}
\par\noindent\begin{minipage}[t]{.49\linewidth}
\vspace{-4mm}
\ExampleInformationTable{
\multicolumn{3}{|l|}{\textbf{Population parameters}}\\ \hline
$\mu$&Population mean&The average duration of all cycling trips\\ \hline
$\mu_{\rm sim}$&Simulation population mean&41 minutes\\ \hline
$\sigma$&Population standard deviation&14 minutes\\ \hline
\multicolumn{3}{|l|}{\textbf{Sample statistics}}\\ \hline
$\bar{x}$&Sample mean&40 minutes\\ \hline
$n$&Sample size&49 trips\\ \hline
$\mathrm{SE}$&Standard error&2 minutes\\ \hline}
\end{minipage}\hfill
\begin{minipage}[t]{.49\linewidth}
\vspace{0pt}
\ExampleInformationTable{
\multicolumn{3}{|l|}{\textbf{Procedure Settings}}\\ \hline
$C$&Confidence level&95\%\\ \hline
$\alpha$&Significance level&0.05\\ \hline
$\alpha/2$&Tail area&0.025\\ \hline
$z^*$&Critical value&1.960\\ \hline}
\vspace{4mm}
Calculate the standard error: $\mathrm{SE}=\sigma/\sqrt n=14/\sqrt{49}=2$ minutes. This is the spread of means across repeated independent samples of 49 trips.
\end{minipage}\par

\vspace{-3mm}
\subsubsection*{Task 2 solution}

\GuidedEstimationTable{
$\bar{x}$ & Point estimate & \vspace{-2mm} $\bar{x}=\dfrac{\sum x_i}{n}$ (given) 40 minutes: Mean cycling-trip duration\\[14pt] \hline
}
\par\smallskip
\GuidedEstimationTable{
\multicolumn{3}{|c|}{\textbf{95\% confidence}}\\ \hline
$\mathrm{MoE}$ & Margin of error & $\mathrm{MoE}=z^*\mathrm{SE}=1.960(2)=3.92$ minutes \qquad  Uncertainty in mean duration\\[4pt] \hline
$\mathrm{CI}$ & Confidence interval & $L=40-3.92=36.08$ minutes \par $U=40+3.92=43.92$ minutes\par
$\mathrm{CI}=[L,U]=[\bar{x}-\mathrm{MoE},\bar{x}+\mathrm{MoE}] $ \par \qquad \qquad $=[40-3.92,40+3.92]$\par
 $\mathrm{CI}=[36.08,43.92]$ minutes\par Plausible mean durations\\[6pt] \hline
$W$ & Full width & $W=U-L=2\mathrm{MoE}$\par $=43.92-36.08=2(3.92)=7.84$ minutes\par Precision of mean-duration estimate\\[6pt] \hline
}
We are 95\% confident that the population mean cycling-trip duration is in the range of plausible values under the model.

\vspace{-3mm}
\subsubsection*{Task 3 solution}
\vspace{-3mm}
The fixed simulation mean is 41 minutes. Samples 1, 2, 4, and 5 cover it because their endpoints satisfy $L\leq41\leq U$. Sample 3 misses: its lower endpoint, 41.08, exceeds 41.
\par 95\%: $N_{\rm cover}(95\%)=4$; $N(95\%)=5$; $\mathrm{Cov}(95\%)=\dfrac45\times100\%=80\%$.
\par $E_{\rm abs}(95\%)=|\mathrm{Cov}(95\%)-95|=|80-95|=15$ percentage points.

Four of five subsequent intervals cover $\mu_{\mathrm{known}}=41$ minutes, the later established population mean: 80\%, 15 percentage points below nominal 95\%. This finite-sample discrepancy is consistent with the model. Nominal confidence describes the initial procedure's long-run interval coverage; observed coverage need not match it or approach it monotonically.

\clearpage
\stageheading{Guided Example}
\vspace{-4mm}
\problemtitle{2}{Two reports on weekly printing costs}
An initial study of all school offices in a district during a specified week used a random sample of 100 offices, with an average weekly printing cost of 86 thousand COP, to estimate the then unknown district average. The known population standard deviation was 32 thousand COP. Reports used 90\% and 99\% confidence.

Two hundred subsequent independent studies each sampled 100 offices from that same population for that week, using the same standard deviation and both confidence procedures. A later comprehensive population investigation established an average weekly cost of 88 thousand COP. Retrospectively, the 90\% intervals missed this value 23 times and the 99\% intervals missed it 5 times.


\task{1}{Identify and organize the information; calculate SE.}
\par\noindent\begin{minipage}[t]{.49\linewidth}
\vspace{0pt}
\ExampleInformationTable{
\multicolumn{3}{|l|}{\textbf{Population parameters}}\\ \hline
$\mu$&Population mean&District mean weekly printing cost:\par\\[10pt] \hline
$\mu_{\mathrm{known}}$&Later established population mean&\\[16pt] \hline
$\sigma$&Population standard deviation&\\[16pt] \hline
\multicolumn{3}{|l|}{\textbf{Sample statistics}}\\ \hline
$\bar{x}$&Sample mean&\\[10pt] \hline
$n$&Sample size&\\[10pt] \hline
$\mathrm{SE}$&Standard error&\\[10pt] \hline}
\end{minipage}\hfill
\begin{minipage}[t]{.49\linewidth}
\vspace{0pt}
\ExampleInformationTable{
\multicolumn{3}{|l|}{\textbf{Procedure Settings}}\\ \hline
$C$&Confidence level&90\%: \par 99\%: \\[16pt] \hline
$\alpha$&Significance level&90\%: \par 99\%: \\[16pt] \hline
$\alpha/2$&Tail area&90\%: \par 99\%: \\[16pt] \hline
$z^*$&Critical value&90\%: \par 99\%: \\[16pt] \hline}
\end{minipage}\par
Calculate $\mathrm{SE}=\sigma/\sqrt n=32/\sqrt{100}=$ \rule{.7in}{.3pt}. Units and meaning: \rule{1.8in}{.3pt}.

\Needspace{20\baselineskip}
\task{2}{Calculate MoE and state the interval.}
\GuidedEstimationTable{
$\bar{x}$ & Shared point estimate & \\[10pt] \hline
}
\par\smallskip
\noindent\begin{minipage}[t]{.49\linewidth}
\vspace{0pt}
\GuidedConfidenceTable{90\%}
\end{minipage}\hfill
\begin{minipage}[t]{.49\linewidth}
\vspace{0pt}
\GuidedConfidenceTable{99\%}
\end{minipage}\par

\problemtitle{2}{Weekly printing costs -- coverage}
\task{3}{Compare observed interval coverage with nominal confidence.}
\par For each level separately, use $\mu_{\mathrm{known}}=88$ thousand COP to find observed interval coverage and its absolute difference from nominal confidence. Interpret both discrepancies in context.
\par\noindent\begin{minipage}[t]{0.70\linewidth}
\vspace{0pt}
\par 90\%: $N_{\rm cover}=200-23=$ \rule{.5in}{.3pt}; $N=$ \rule{.4in}{.3pt}; $\mathrm{Cov}=$ \rule{1in}{.3pt}.
\par $E_{\rm abs}=|\mathrm{Cov}-90|=$ \rule{.7in}{.3pt} percentage points. 
\par 99\%: $N_{\rm cover}=200-5=$ \rule{.5in}{.3pt}; $N=$ \rule{.4in}{.3pt}; $\mathrm{Cov}=$ \rule{1in}{.3pt}.
\par $E_{\rm abs}=|\mathrm{Cov}-99|=$ \rule{.7in}{.3pt} percentage points. 
\end{minipage}\hfill
\begin{minipage}[t]{0.28\linewidth}
\vspace{0pt}
\par Explain finite-sample variation under the model. Must observed coverage equal nominal confidence or approach it monotonically as studies are added?
\end{minipage}\par


\clearpage
\stageheading{Independent Practice}
Complete the following three tasks for each problem (3--6).

\task{1}{Identify and organize the information; calculate SE.}
\task{2}{Calculate MoE and state the interval.}
\task{3}{Compare observed interval coverage with nominal confidence.}

\problemtitle{3}{Repeated estimates of locker rental charges}
\noindent\begin{minipage}[t]{0.76\linewidth}
\vspace{0pt}
An initial study of all customers at a sports centre during a specified week used a random sample of 64 customers, with an average weekly locker rental charge of 72 thousand COP, to estimate the then unknown population average. The known population standard deviation was 20 thousand COP, and the report used 95\% confidence.

Five subsequent independent studies each sampled 64 customers from that same population for that week, using the same standard deviation and 95\% procedure. A later comprehensive population investigation established an average weekly charge of 70 thousand COP, allowing retrospective assessment. The table gives interval endpoints in thousand COP.
\end{minipage}\hfill
\begin{minipage}[t]{0.22\linewidth}
\vspace{0pt}
\centering
\small
\setlength{\tabcolsep}{4pt}
\begin{tabularx}{\linewidth}{@{}l>{\raggedleft\arraybackslash}X>{\raggedleft\arraybackslash}X@{}}
	\toprule
	Study & Lower endpoint & Upper endpoint \\
	\midrule
	1 & 64.1 & 73.9 \\
	2 & 67.1 & 76.9 \\
	3 & 70.1 & 79.9 \\
	4 & 63.1 & 72.9 \\
	5 & 65.1 & 74.9 \\
	\bottomrule
\end{tabularx}
\end{minipage}\par
\problemtitle{4}{Two reports on monthly pet-care spending}
An initial study of all registered households in a district during a specified month used a random sample of 144 households, with average monthly pet-care spending of 145 thousand COP, to estimate the then unknown district average. The known population standard deviation was 42 thousand COP. Reports used 90\% and 99\% confidence.

One hundred fifty subsequent independent studies each sampled 144 households from that same population for that month, using the same standard deviation and both confidence procedures. A later comprehensive population investigation established average monthly spending of 148 thousand COP. Retrospectively, the 90\% intervals missed this value 17 times and the 99\% intervals missed it 3 times.

\problemtitle{5}{Studies of clinic waiting times}
\noindent\begin{minipage}[t]{0.83\linewidth}
\vspace{0pt}
An initial study of all appointments at a clinic during a specified month used a random sample of 49 appointments, with an average waiting time of 26 minutes, to estimate the then unknown monthly population average. The known population standard deviation was 14 minutes, and the report used 90\% confidence.

Twenty subsequent independent studies each sampled 49 appointments from that same month's population, using the same standard deviation and 90\% procedure. A later comprehensive population investigation established an average waiting time of 27 minutes. The table groups counts of intervals covering this later established mean by study order.
\end{minipage}\hfill
\begin{minipage}[t]{0.15\linewidth}
\vspace{0pt}
\centering
\small
\setlength{\tabcolsep}{4pt}
\begin{tabularx}{\linewidth}{@{}>{\raggedright\arraybackslash}X>{\raggedleft\arraybackslash}X@{}}\toprule
Studies in order & Covering intervals\\\midrule
1--5 & 4\\
6--10 & 5\\
11--15 & 4\\
16--20 & 3\\
\bottomrule\end{tabularx}
\end{minipage}\par

\problemtitle{6}{Three confidence levels for library processing times}
\noindent\begin{minipage}[t]{0.70\linewidth}
\vspace{0pt}
An initial study of all books returned to a library during a specified month used a random sample of 100 books, with an average processing time of 24 minutes, to estimate the then unknown monthly population average. The known population standard deviation was 15 minutes. Reports used 90\%, 95\%, and 99\% confidence.

Two hundred fifty subsequent independent studies each sampled 100 books from that same month's population, using the same standard deviation and all three confidence procedures. A later comprehensive population investigation established an average processing time of 25 minutes. Within each study, intervals share its own sample average and are nested from narrowest at 90\% to widest at 99\%. The covering categories are mutually exclusive; each study appears once.
\end{minipage}\hfill
\begin{minipage}[t]{0.28\linewidth}
\vspace{0pt}
\centering
\small
\setlength{\tabcolsep}{4pt}
\begin{tabularx}{\linewidth}{@{}>{\raggedright\arraybackslash}Xr@{}}\toprule
Intervals covering the later established mean & Studies\\\midrule
All three & 219\\
Only 95\% and 99\% & 18\\
Only 99\% & 10\\
None & 3\\
\bottomrule\end{tabularx}
\end{minipage}\par


% Appended review only: the original classroom body above is preserved verbatim.
\clearpage
\newgeometry{margin=0.65in,top=0.55in,bottom=0.65in}
\begingroup
\linespread{1.12}\selectfont
\setlength{\parskip}{4pt}
\CtwoPracticeIntervalPalette
\section*{Solutions}
Solutions to problems 2--6. Problem 1 already has its worked solution in the activity.
The highlighted names, symbols, substitutions, and results use the practice document's semantic roles.
In every problem, $\CtwoMean{\mu_{\mathrm{known}}}$ is the later established value of the same fixed population mean $\CtwoMean{\mu}$.
\subsection*{2. Two reports on weekly printing costs}
\subsubsection*{Task 1 solution. Organize the information and calculate SE}
\SolutionInformationTable{
$\CtwoMean{\mu}$ & \CtwoMean{Population mean} & Population parameters & The district mean weekly printing cost for all school offices in the specified week, in thousand COP; fixed and initially unknown. \\[6pt]
$\CtwoMean{\mu_{\mathrm{known}}}$ & \CtwoMean{Later established} \CtwoMean{population mean} &  & $\CtwoMean{88}$ thousand COP; established by the later comprehensive population investigation. \\[6pt]
$\CtwoSD{\sigma}$ & \CtwoSD{Population} \CtwoSD{standard deviation} &  & $\CtwoSD{32}$ thousand COP; known spread of individual observations. \\[6pt]
\midrule
$\CtwoPoint{\bar{x}}$ & \CtwoPoint{Sample mean} & Sample statistics & $\CtwoPoint{86}$ thousand COP; the initial point estimate of $\CtwoMean{\mu}$. \\[6pt]
$\CtwoSize{n}$ & \CtwoSize{Sample size} &  & $\CtwoSize{100}$ offices in the initial sample and each subsequent study. \\[6pt]
$\CtwoSE{\mathrm{SE}}$ & \CtwoSE{Standard error} &  & $\CtwoSE{\mathrm{SE}}=\CtwoSD{\sigma}/\sqrt{\CtwoSize{n}}$\par$=\CtwoSD{32}/\sqrt{\CtwoSize{100}}=\CtwoSE{3.2}$ thousand COP. Spread of means across repeated independent samples of 100 offices. \\[6pt]
\midrule
$\CtwoConfidence{C}$ & \CtwoConfidence{Confidence} \CtwoConfidence{level} & Procedure Settings & $\CtwoConfidence{90\%}$: $\CtwoConfidence{C}=\CtwoConfidence{0.90}$\par $\CtwoConfidence{99\%}$: $\CtwoConfidence{C}=\CtwoConfidence{0.99}$. \\[6pt]
$\CtwoSignificance{\alpha}$ & \CtwoSignificance{Significance} \CtwoSignificance{level} &  & $\CtwoConfidence{90\%}$: $\CtwoSignificance{\alpha}=1-\CtwoConfidence{0.90}=\CtwoSignificance{0.10}$\par $\CtwoConfidence{99\%}$: $\CtwoSignificance{\alpha}=1-\CtwoConfidence{0.99}=\CtwoSignificance{0.01}$. \\[6pt]
$\CtwoSignificance{\alpha/2}$ & \CtwoSignificance{Tail} \CtwoSignificance{area} &  & $\CtwoConfidence{90\%}$: $\CtwoSignificance{\alpha/2}=\CtwoSignificance{0.10}/2=\CtwoSignificance{0.05}$\par $\CtwoConfidence{99\%}$: $\CtwoSignificance{\alpha/2}=\CtwoSignificance{0.01}/2=\CtwoSignificance{0.005}$. \\[6pt]
$\CtwoCritical{z^*}$ & \CtwoCritical{Critical} \CtwoCritical{value} &  & $\CtwoConfidence{90\%}$: $\CtwoCritical{z^*}=\CtwoCritical{1.645}$\par $\CtwoConfidence{99\%}$: $\CtwoCritical{z^*}=\CtwoCritical{2.576}$. \\[6pt]
}
The observations are independent and random, $\CtwoSD{\sigma}$ is known, and the sample mean is normally distributed under the stated model. No sample SD is required.
\clearpage
\subsubsection*{Task 2 solution. Calculate MoE and state the interval}
\textbf{$\CtwoConfidence{90\%}$ confidence}
\SolutionTaskTwoTable{
$\CtwoPoint{\bar{x}}$ & \CtwoPoint{Point estimate} & $\CtwoPoint{\bar{x}}=\sum x_i/\CtwoSize{n}$\par Given sample mean & $\CtwoPoint{86}$\par thousand COP & Estimates the population mean. \\[6pt]
\midrule
$\CtwoMoE{\mathrm{MoE}}$ & \CtwoMoE{Margin of error} & $\CtwoMoE{\mathrm{MoE}}=\CtwoCritical{z^*}\times\CtwoSE{\mathrm{SE}}$\par$=\CtwoCritical{1.645}\times\CtwoSE{3.2}$ & $\CtwoMoE{5.264}$\par thousand COP & Distance from centre to either bound. \\[6pt]
$\CtwoEndpoint{L}$ & \CtwoEndpoint{Lower} \CtwoEndpoint{endpoint} & $\CtwoEndpoint{L}=\CtwoPoint{\bar{x}}-\CtwoMoE{\mathrm{MoE}}$\par$=\CtwoPoint{86}-\CtwoMoE{5.264}$ & $\CtwoEndpoint{80.736}$\par thousand COP & Lower bound for the population mean. \\[6pt]
$\CtwoEndpoint{U}$ & \CtwoEndpoint{Upper} \CtwoEndpoint{endpoint} & $\CtwoEndpoint{U}=\CtwoPoint{\bar{x}}+\CtwoMoE{\mathrm{MoE}}$\par$=\CtwoPoint{86}+\CtwoMoE{5.264}$ & $\CtwoEndpoint{91.264}$\par thousand COP & Upper bound for the population mean. \\[6pt]
$\CtwoCI{\mathrm{CI}}$ & \CtwoCI{Confidence} \CtwoCI{interval} & $\CtwoCI{\mathrm{CI}}=[\CtwoEndpoint{L},\CtwoEndpoint{U}]$\par$=[\CtwoPoint{86}-\CtwoMoE{5.264},$\par$\CtwoPoint{86}+\CtwoMoE{5.264}]$ & $[\CtwoEndpoint{80.736},$\par$\CtwoEndpoint{91.264}]$\par thousand COP & Plausible values for the population mean under the model. \\[6pt]
$\CtwoWidth{W}$ & \CtwoWidth{Full width} & $\CtwoWidth{W}=\CtwoEndpoint{U}-\CtwoEndpoint{L}$\par$=\CtwoEndpoint{91.264}-\CtwoEndpoint{80.736}$\par$=2\CtwoMoE{\mathrm{MoE}}=2(\CtwoMoE{5.264})$ & $\CtwoWidth{10.528}$\par thousand COP & Total span; smaller width means greater precision. \\[6pt]
}
We are \CtwoConfidence{90\%} confident that the district mean weekly printing cost for all school offices in the specified week lies between $\CtwoEndpoint{80.736}$ and $\CtwoEndpoint{91.264}$ thousand COP. Under the model, \CtwoConfidence{90\%} of intervals constructed by this procedure would cover the fixed $\CtwoMean{\mu}$ in the long run. This is not a probability assigned to the fixed mean after calculating this interval.
\clearpage
\subsubsection*{Task 2 solution (continued). Calculate MoE and state the interval}
\textbf{$\CtwoConfidence{99\%}$ confidence}
\SolutionTaskTwoTable{
\multicolumn{5}{@{}l}{Shared \CtwoPoint{point estimate}: $\CtwoPoint{\bar{x}}$=$\CtwoPoint{86}$ thousand COP.}\\[6pt]
\midrule
$\CtwoMoE{\mathrm{MoE}}$ & \CtwoMoE{Margin of error} & $\CtwoMoE{\mathrm{MoE}}=\CtwoCritical{z^*}\times\CtwoSE{\mathrm{SE}}$\par$=\CtwoCritical{2.576}\times\CtwoSE{3.2}$ & $\CtwoMoE{8.2432}$\par thousand COP & Distance from centre to either bound. \\[6pt]
$\CtwoEndpoint{L}$ & \CtwoEndpoint{Lower} \CtwoEndpoint{endpoint} & $\CtwoEndpoint{L}=\CtwoPoint{\bar{x}}-\CtwoMoE{\mathrm{MoE}}$\par$=\CtwoPoint{86}-\CtwoMoE{8.2432}$ & $\CtwoEndpoint{77.7568}$\par thousand COP & Lower bound for the population mean. \\[6pt]
$\CtwoEndpoint{U}$ & \CtwoEndpoint{Upper} \CtwoEndpoint{endpoint} & $\CtwoEndpoint{U}=\CtwoPoint{\bar{x}}+\CtwoMoE{\mathrm{MoE}}$\par$=\CtwoPoint{86}+\CtwoMoE{8.2432}$ & $\CtwoEndpoint{94.2432}$\par thousand COP & Upper bound for the population mean. \\[6pt]
$\CtwoCI{\mathrm{CI}}$ & \CtwoCI{Confidence} \CtwoCI{interval} & $\CtwoCI{\mathrm{CI}}=[\CtwoEndpoint{L},\CtwoEndpoint{U}]$\par$=[\CtwoPoint{86}-\CtwoMoE{8.2432},$\par$\CtwoPoint{86}+\CtwoMoE{8.2432}]$ & $[\CtwoEndpoint{77.7568},$\par$\CtwoEndpoint{94.2432}]$\par thousand COP & Plausible values for the population mean under the model. \\[6pt]
$\CtwoWidth{W}$ & \CtwoWidth{Full width} & $\CtwoWidth{W}=\CtwoEndpoint{U}-\CtwoEndpoint{L}$\par$=\CtwoEndpoint{94.2432}-\CtwoEndpoint{77.7568}$\par$=2\CtwoMoE{\mathrm{MoE}}=2(\CtwoMoE{8.2432})$ & $\CtwoWidth{16.4864}$\par thousand COP & Total span; smaller width means greater precision. \\[6pt]
}
We are \CtwoConfidence{99\%} confident that the district mean weekly printing cost for all school offices in the specified week lies between $\CtwoEndpoint{77.7568}$ and $\CtwoEndpoint{94.2432}$ thousand COP. Under the model, \CtwoConfidence{99\%} of intervals constructed by this procedure would cover the fixed $\CtwoMean{\mu}$ in the long run. This is not a probability assigned to the fixed mean after calculating this interval.
\clearpage
\subsubsection*{Task 3 solution. Calculate coverage, absolute error, and interpret confidence}
For these subsequent studies, an interval covers exactly when $\CtwoEndpoint{L_i}\leq\CtwoMean{\mu_{\mathrm{known}}}\leq\CtwoEndpoint{U_i}$, including equality. The later established mean is $\CtwoMean{88}$ thousand COP. Each interval is centred on its own study sample mean.
\Needspace{10\baselineskip}
\textbf{$\CtwoConfidence{90\%}$ confidence}
\begin{align*}
\CtwoCoverCount(\CtwoConfidence{90\%})&=\CtwoCoverage{200}-\CtwoCoverage{23}=\CtwoCoverage{177},\qquad \CtwoTotalCount(\CtwoConfidence{90\%})=\CtwoCoverage{200};\\
\CtwoCoverage{\mathrm{Cov}}(\CtwoConfidence{90\%})&=\frac{\CtwoCoverage{177}}{\CtwoCoverage{200}}\times100\%\approx\CtwoCoverage{88.50\%};\\
\CtwoError{E_{\rm abs}}(\CtwoConfidence{90\%})&=\left|\frac{\CtwoCoverage{17700}}{\CtwoCoverage{200}}-\CtwoConfidence{90}\right|\approx\CtwoError{1.50}\ \text{percentage points}.
\end{align*}
Here $\CtwoCoverage{177}$ of $\CtwoCoverage{200}$ subsequent intervals cover the district mean weekly printing cost for all school offices in the specified week: observed coverage is below nominal \CtwoConfidence{90\%} confidence by \CtwoError{1.50} \CtwoError{percentage points}.
\Needspace{10\baselineskip}
\textbf{$\CtwoConfidence{99\%}$ confidence}
\begin{align*}
\CtwoCoverCount(\CtwoConfidence{99\%})&=\CtwoCoverage{200}-\CtwoCoverage{5}=\CtwoCoverage{195},\qquad \CtwoTotalCount(\CtwoConfidence{99\%})=\CtwoCoverage{200};\\
\CtwoCoverage{\mathrm{Cov}}(\CtwoConfidence{99\%})&=\frac{\CtwoCoverage{195}}{\CtwoCoverage{200}}\times100\%\approx\CtwoCoverage{97.50\%};\\
\CtwoError{E_{\rm abs}}(\CtwoConfidence{99\%})&=\left|\frac{\CtwoCoverage{19500}}{\CtwoCoverage{200}}-\CtwoConfidence{99}\right|\approx\CtwoError{1.50}\ \text{percentage points}.
\end{align*}
Here $\CtwoCoverage{195}$ of $\CtwoCoverage{200}$ subsequent intervals cover the district mean weekly printing cost for all school offices in the specified week: observed coverage is below nominal \CtwoConfidence{99\%} confidence by \CtwoError{1.50} \CtwoError{percentage points}.
For each study the intervals share its sample mean and \CtwoSE{SE}. Increasing \CtwoConfidence{confidence} increases $\CtwoCritical{z^*}$, \CtwoMoE{MoE}, and \CtwoWidth{full width}. The narrower interval lies inside each wider partner, so covering counts cannot decrease with confidence on these same studies. 
These are retrospective studies of the same fixed population and target period. Finite-run \CtwoCoverage{observed coverage} need not equal \CtwoConfidence{nominal confidence} or approach it monotonically as studies are added. Under the model, confidence describes long-run coverage of the sampling procedure, including the one used for the initial report. Each realized interval either covers or misses the fixed mean; coverage is not the fraction of individual observations inside the initial interval. A smaller finite-run \CtwoError{absolute error} does not establish superior long-run performance.
\clearpage
\subsection*{3. Repeated estimates of locker rental charges}
\subsubsection*{Task 1 solution. Organize the information and calculate SE}
\SolutionInformationTable{
$\CtwoMean{\mu}$ & \CtwoMean{Population mean} & Population parameters & The mean weekly locker rental charge for all sports-centre customers in the specified week, in thousand COP; fixed and initially unknown. \\[6pt]
$\CtwoMean{\mu_{\mathrm{known}}}$ & \CtwoMean{Later established} \CtwoMean{population mean} &  & $\CtwoMean{70}$ thousand COP; established by the later comprehensive population investigation. \\[6pt]
$\CtwoSD{\sigma}$ & \CtwoSD{Population} \CtwoSD{standard deviation} &  & $\CtwoSD{20}$ thousand COP; known spread of individual observations. \\[6pt]
\midrule
$\CtwoPoint{\bar{x}}$ & \CtwoPoint{Sample mean} & Sample statistics & $\CtwoPoint{72}$ thousand COP; the initial point estimate of $\CtwoMean{\mu}$. \\[6pt]
$\CtwoSize{n}$ & \CtwoSize{Sample size} &  & $\CtwoSize{64}$ customers in the initial sample and each subsequent study. \\[6pt]
$\CtwoSE{\mathrm{SE}}$ & \CtwoSE{Standard error} &  & $\CtwoSE{\mathrm{SE}}=\CtwoSD{\sigma}/\sqrt{\CtwoSize{n}}$\par$=\CtwoSD{20}/\sqrt{\CtwoSize{64}}=\CtwoSE{2.5}$ thousand COP. Spread of means across repeated independent samples of 64 customers. \\[6pt]
\midrule
$\CtwoConfidence{C}$ & \CtwoConfidence{Confidence} \CtwoConfidence{level} & Procedure Settings & $\CtwoConfidence{95\%}$: $\CtwoConfidence{C}=\CtwoConfidence{0.95}$. \\[6pt]
$\CtwoSignificance{\alpha}$ & \CtwoSignificance{Significance} \CtwoSignificance{level} &  & $\CtwoConfidence{95\%}$: $\CtwoSignificance{\alpha}=1-\CtwoConfidence{0.95}=\CtwoSignificance{0.05}$. \\[6pt]
$\CtwoSignificance{\alpha/2}$ & \CtwoSignificance{Tail} \CtwoSignificance{area} &  & $\CtwoConfidence{95\%}$: $\CtwoSignificance{\alpha/2}=\CtwoSignificance{0.05}/2=\CtwoSignificance{0.025}$. \\[6pt]
$\CtwoCritical{z^*}$ & \CtwoCritical{Critical} \CtwoCritical{value} &  & $\CtwoConfidence{95\%}$: $\CtwoCritical{z^*}=\CtwoCritical{1.960}$. \\[6pt]
}
The observations are independent and random, $\CtwoSD{\sigma}$ is known, and the sample mean is normally distributed under the stated model. No sample SD is required.
\clearpage
\subsubsection*{Task 2 solution. Calculate MoE and state the interval}
\textbf{$\CtwoConfidence{95\%}$ confidence}
\SolutionTaskTwoTable{
$\CtwoPoint{\bar{x}}$ & \CtwoPoint{Point estimate} & $\CtwoPoint{\bar{x}}=\sum x_i/\CtwoSize{n}$\par Given sample mean & $\CtwoPoint{72}$\par thousand COP & Estimates the population mean. \\[6pt]
\midrule
$\CtwoMoE{\mathrm{MoE}}$ & \CtwoMoE{Margin of error} & $\CtwoMoE{\mathrm{MoE}}=\CtwoCritical{z^*}\times\CtwoSE{\mathrm{SE}}$\par$=\CtwoCritical{1.960}\times\CtwoSE{2.5}$ & $\CtwoMoE{4.9}$\par thousand COP & Distance from centre to either bound. \\[6pt]
$\CtwoEndpoint{L}$ & \CtwoEndpoint{Lower} \CtwoEndpoint{endpoint} & $\CtwoEndpoint{L}=\CtwoPoint{\bar{x}}-\CtwoMoE{\mathrm{MoE}}$\par$=\CtwoPoint{72}-\CtwoMoE{4.9}$ & $\CtwoEndpoint{67.1}$\par thousand COP & Lower bound for the population mean. \\[6pt]
$\CtwoEndpoint{U}$ & \CtwoEndpoint{Upper} \CtwoEndpoint{endpoint} & $\CtwoEndpoint{U}=\CtwoPoint{\bar{x}}+\CtwoMoE{\mathrm{MoE}}$\par$=\CtwoPoint{72}+\CtwoMoE{4.9}$ & $\CtwoEndpoint{76.9}$\par thousand COP & Upper bound for the population mean. \\[6pt]
$\CtwoCI{\mathrm{CI}}$ & \CtwoCI{Confidence} \CtwoCI{interval} & $\CtwoCI{\mathrm{CI}}=[\CtwoEndpoint{L},\CtwoEndpoint{U}]$\par$=[\CtwoPoint{72}-\CtwoMoE{4.9},$\par$\CtwoPoint{72}+\CtwoMoE{4.9}]$ & $[\CtwoEndpoint{67.1},$\par$\CtwoEndpoint{76.9}]$\par thousand COP & Plausible values for the population mean under the model. \\[6pt]
$\CtwoWidth{W}$ & \CtwoWidth{Full width} & $\CtwoWidth{W}=\CtwoEndpoint{U}-\CtwoEndpoint{L}$\par$=\CtwoEndpoint{76.9}-\CtwoEndpoint{67.1}$\par$=2\CtwoMoE{\mathrm{MoE}}=2(\CtwoMoE{4.9})$ & $\CtwoWidth{9.8}$\par thousand COP & Total span; smaller width means greater precision. \\[6pt]
}
We are \CtwoConfidence{95\%} confident that the mean weekly locker rental charge for all sports-centre customers in the specified week lies between $\CtwoEndpoint{67.1}$ and $\CtwoEndpoint{76.9}$ thousand COP. Under the model, \CtwoConfidence{95\%} of intervals constructed by this procedure would cover the fixed $\CtwoMean{\mu}$ in the long run. This is not a probability assigned to the fixed mean after calculating this interval.
\clearpage
\subsubsection*{Task 3 solution. Calculate coverage, absolute error, and interpret confidence}
For these subsequent studies, an interval covers exactly when $\CtwoEndpoint{L_i}\leq\CtwoMean{\mu_{\mathrm{known}}}\leq\CtwoEndpoint{U_i}$, including equality. The later established mean is $\CtwoMean{70}$ thousand COP. Each interval is centred on its own study sample mean.
Studies 1, 2, 4, and 5 cover: $\CtwoEndpoint{64.1}\leq\CtwoMean{70}\leq\CtwoEndpoint{73.9}$, $\CtwoEndpoint{67.1}\leq\CtwoMean{70}\leq\CtwoEndpoint{76.9}$, $\CtwoEndpoint{63.1}\leq\CtwoMean{70}\leq\CtwoEndpoint{72.9}$, and $\CtwoEndpoint{65.1}\leq\CtwoMean{70}\leq\CtwoEndpoint{74.9}$. Study 3 misses because its \CtwoEndpoint{lower endpoint} $\CtwoEndpoint{70.1}$ exceeds $\CtwoMean{70}$.
\Needspace{10\baselineskip}
\textbf{$\CtwoConfidence{95\%}$ confidence}
\begin{align*}
\CtwoCoverCount(\CtwoConfidence{95\%})&=\CtwoCoverage{4}=\CtwoCoverage{4},\qquad \CtwoTotalCount(\CtwoConfidence{95\%})=\CtwoCoverage{5};\\
\CtwoCoverage{\mathrm{Cov}}(\CtwoConfidence{95\%})&=\frac{\CtwoCoverage{4}}{\CtwoCoverage{5}}\times100\%=\CtwoCoverage{80\%};\\
\CtwoError{E_{\rm abs}}(\CtwoConfidence{95\%})&=\left|\CtwoCoverage{80}-\CtwoConfidence{95}\right|=\CtwoError{15}\ \text{percentage points}.
\end{align*}
Here $\CtwoCoverage{4}$ of $\CtwoCoverage{5}$ subsequent intervals cover the mean weekly locker rental charge for all sports-centre customers in the specified week: observed coverage is below nominal \CtwoConfidence{95\%} confidence by \CtwoError{15} \CtwoError{percentage points}.
These are retrospective studies of the same fixed population and target period. Finite-run \CtwoCoverage{observed coverage} need not equal \CtwoConfidence{nominal confidence} or approach it monotonically as studies are added. Under the model, confidence describes long-run coverage of the sampling procedure, including the one used for the initial report. Each realized interval either covers or misses the fixed mean; coverage is not the fraction of individual observations inside the initial interval. A smaller finite-run \CtwoError{absolute error} does not establish superior long-run performance.
\clearpage
\subsection*{4. Two reports on monthly pet-care spending}
\subsubsection*{Task 1 solution. Organize the information and calculate SE}
\SolutionInformationTable{
$\CtwoMean{\mu}$ & \CtwoMean{Population mean} & Population parameters & The district mean monthly pet-care spending for all registered households in the specified month, in thousand COP; fixed and initially unknown. \\[6pt]
$\CtwoMean{\mu_{\mathrm{known}}}$ & \CtwoMean{Later established} \CtwoMean{population mean} &  & $\CtwoMean{148}$ thousand COP; established by the later comprehensive population investigation. \\[6pt]
$\CtwoSD{\sigma}$ & \CtwoSD{Population} \CtwoSD{standard deviation} &  & $\CtwoSD{42}$ thousand COP; known spread of individual observations. \\[6pt]
\midrule
$\CtwoPoint{\bar{x}}$ & \CtwoPoint{Sample mean} & Sample statistics & $\CtwoPoint{145}$ thousand COP; the initial point estimate of $\CtwoMean{\mu}$. \\[6pt]
$\CtwoSize{n}$ & \CtwoSize{Sample size} &  & $\CtwoSize{144}$ households in the initial sample and each subsequent study. \\[6pt]
$\CtwoSE{\mathrm{SE}}$ & \CtwoSE{Standard error} &  & $\CtwoSE{\mathrm{SE}}=\CtwoSD{\sigma}/\sqrt{\CtwoSize{n}}$\par$=\CtwoSD{42}/\sqrt{\CtwoSize{144}}=\CtwoSE{3.5}$ thousand COP. Spread of means across repeated independent samples of 144 households. \\[6pt]
\midrule
$\CtwoConfidence{C}$ & \CtwoConfidence{Confidence} \CtwoConfidence{level} & Procedure Settings & $\CtwoConfidence{90\%}$: $\CtwoConfidence{C}=\CtwoConfidence{0.90}$\par $\CtwoConfidence{99\%}$: $\CtwoConfidence{C}=\CtwoConfidence{0.99}$. \\[6pt]
$\CtwoSignificance{\alpha}$ & \CtwoSignificance{Significance} \CtwoSignificance{level} &  & $\CtwoConfidence{90\%}$: $\CtwoSignificance{\alpha}=1-\CtwoConfidence{0.90}=\CtwoSignificance{0.10}$\par $\CtwoConfidence{99\%}$: $\CtwoSignificance{\alpha}=1-\CtwoConfidence{0.99}=\CtwoSignificance{0.01}$. \\[6pt]
$\CtwoSignificance{\alpha/2}$ & \CtwoSignificance{Tail} \CtwoSignificance{area} &  & $\CtwoConfidence{90\%}$: $\CtwoSignificance{\alpha/2}=\CtwoSignificance{0.10}/2=\CtwoSignificance{0.05}$\par $\CtwoConfidence{99\%}$: $\CtwoSignificance{\alpha/2}=\CtwoSignificance{0.01}/2=\CtwoSignificance{0.005}$. \\[6pt]
$\CtwoCritical{z^*}$ & \CtwoCritical{Critical} \CtwoCritical{value} &  & $\CtwoConfidence{90\%}$: $\CtwoCritical{z^*}=\CtwoCritical{1.645}$\par $\CtwoConfidence{99\%}$: $\CtwoCritical{z^*}=\CtwoCritical{2.576}$. \\[6pt]
}
The observations are independent and random, $\CtwoSD{\sigma}$ is known, and the sample mean is normally distributed under the stated model. No sample SD is required.
\clearpage
\subsubsection*{Task 2 solution. Calculate MoE and state the interval}
\textbf{$\CtwoConfidence{90\%}$ confidence}
\SolutionTaskTwoTable{
$\CtwoPoint{\bar{x}}$ & \CtwoPoint{Point estimate} & $\CtwoPoint{\bar{x}}=\sum x_i/\CtwoSize{n}$\par Given sample mean & $\CtwoPoint{145}$\par thousand COP & Estimates the population mean. \\[6pt]
\midrule
$\CtwoMoE{\mathrm{MoE}}$ & \CtwoMoE{Margin of error} & $\CtwoMoE{\mathrm{MoE}}=\CtwoCritical{z^*}\times\CtwoSE{\mathrm{SE}}$\par$=\CtwoCritical{1.645}\times\CtwoSE{3.5}$ & $\CtwoMoE{5.7575}$\par thousand COP & Distance from centre to either bound. \\[6pt]
$\CtwoEndpoint{L}$ & \CtwoEndpoint{Lower} \CtwoEndpoint{endpoint} & $\CtwoEndpoint{L}=\CtwoPoint{\bar{x}}-\CtwoMoE{\mathrm{MoE}}$\par$=\CtwoPoint{145}-\CtwoMoE{5.7575}$ & $\CtwoEndpoint{139.2425}$\par thousand COP & Lower bound for the population mean. \\[6pt]
$\CtwoEndpoint{U}$ & \CtwoEndpoint{Upper} \CtwoEndpoint{endpoint} & $\CtwoEndpoint{U}=\CtwoPoint{\bar{x}}+\CtwoMoE{\mathrm{MoE}}$\par$=\CtwoPoint{145}+\CtwoMoE{5.7575}$ & $\CtwoEndpoint{150.7575}$\par thousand COP & Upper bound for the population mean. \\[6pt]
$\CtwoCI{\mathrm{CI}}$ & \CtwoCI{Confidence} \CtwoCI{interval} & $\CtwoCI{\mathrm{CI}}=[\CtwoEndpoint{L},\CtwoEndpoint{U}]$\par$=[\CtwoPoint{145}-\CtwoMoE{5.7575},$\par$\CtwoPoint{145}+\CtwoMoE{5.7575}]$ & $[\CtwoEndpoint{139.2425},$\par$\CtwoEndpoint{150.7575}]$\par thousand COP & Plausible values for the population mean under the model. \\[6pt]
$\CtwoWidth{W}$ & \CtwoWidth{Full width} & $\CtwoWidth{W}=\CtwoEndpoint{U}-\CtwoEndpoint{L}$\par$=\CtwoEndpoint{150.7575}-\CtwoEndpoint{139.2425}$\par$=2\CtwoMoE{\mathrm{MoE}}=2(\CtwoMoE{5.7575})$ & $\CtwoWidth{11.515}$\par thousand COP & Total span; smaller width means greater precision. \\[6pt]
}
We are \CtwoConfidence{90\%} confident that the district mean monthly pet-care spending for all registered households in the specified month lies between $\CtwoEndpoint{139.2425}$ and $\CtwoEndpoint{150.7575}$ thousand COP. Under the model, \CtwoConfidence{90\%} of intervals constructed by this procedure would cover the fixed $\CtwoMean{\mu}$ in the long run. This is not a probability assigned to the fixed mean after calculating this interval.
\clearpage
\subsubsection*{Task 2 solution (continued). Calculate MoE and state the interval}
\textbf{$\CtwoConfidence{99\%}$ confidence}
\SolutionTaskTwoTable{
\multicolumn{5}{@{}l}{Shared \CtwoPoint{point estimate}: $\CtwoPoint{\bar{x}}$=$\CtwoPoint{145}$ thousand COP.}\\[6pt]
\midrule
$\CtwoMoE{\mathrm{MoE}}$ & \CtwoMoE{Margin of error} & $\CtwoMoE{\mathrm{MoE}}=\CtwoCritical{z^*}\times\CtwoSE{\mathrm{SE}}$\par$=\CtwoCritical{2.576}\times\CtwoSE{3.5}$ & $\CtwoMoE{9.016}$\par thousand COP & Distance from centre to either bound. \\[6pt]
$\CtwoEndpoint{L}$ & \CtwoEndpoint{Lower} \CtwoEndpoint{endpoint} & $\CtwoEndpoint{L}=\CtwoPoint{\bar{x}}-\CtwoMoE{\mathrm{MoE}}$\par$=\CtwoPoint{145}-\CtwoMoE{9.016}$ & $\CtwoEndpoint{135.984}$\par thousand COP & Lower bound for the population mean. \\[6pt]
$\CtwoEndpoint{U}$ & \CtwoEndpoint{Upper} \CtwoEndpoint{endpoint} & $\CtwoEndpoint{U}=\CtwoPoint{\bar{x}}+\CtwoMoE{\mathrm{MoE}}$\par$=\CtwoPoint{145}+\CtwoMoE{9.016}$ & $\CtwoEndpoint{154.016}$\par thousand COP & Upper bound for the population mean. \\[6pt]
$\CtwoCI{\mathrm{CI}}$ & \CtwoCI{Confidence} \CtwoCI{interval} & $\CtwoCI{\mathrm{CI}}=[\CtwoEndpoint{L},\CtwoEndpoint{U}]$\par$=[\CtwoPoint{145}-\CtwoMoE{9.016},$\par$\CtwoPoint{145}+\CtwoMoE{9.016}]$ & $[\CtwoEndpoint{135.984},$\par$\CtwoEndpoint{154.016}]$\par thousand COP & Plausible values for the population mean under the model. \\[6pt]
$\CtwoWidth{W}$ & \CtwoWidth{Full width} & $\CtwoWidth{W}=\CtwoEndpoint{U}-\CtwoEndpoint{L}$\par$=\CtwoEndpoint{154.016}-\CtwoEndpoint{135.984}$\par$=2\CtwoMoE{\mathrm{MoE}}=2(\CtwoMoE{9.016})$ & $\CtwoWidth{18.032}$\par thousand COP & Total span; smaller width means greater precision. \\[6pt]
}
We are \CtwoConfidence{99\%} confident that the district mean monthly pet-care spending for all registered households in the specified month lies between $\CtwoEndpoint{135.984}$ and $\CtwoEndpoint{154.016}$ thousand COP. Under the model, \CtwoConfidence{99\%} of intervals constructed by this procedure would cover the fixed $\CtwoMean{\mu}$ in the long run. This is not a probability assigned to the fixed mean after calculating this interval.
\clearpage
\subsubsection*{Task 3 solution. Calculate coverage, absolute error, and interpret confidence}
For these subsequent studies, an interval covers exactly when $\CtwoEndpoint{L_i}\leq\CtwoMean{\mu_{\mathrm{known}}}\leq\CtwoEndpoint{U_i}$, including equality. The later established mean is $\CtwoMean{148}$ thousand COP. Each interval is centred on its own study sample mean.
\Needspace{10\baselineskip}
\textbf{$\CtwoConfidence{90\%}$ confidence}
\begin{align*}
\CtwoCoverCount(\CtwoConfidence{90\%})&=\CtwoCoverage{150}-\CtwoCoverage{17}=\CtwoCoverage{133},\qquad \CtwoTotalCount(\CtwoConfidence{90\%})=\CtwoCoverage{150};\\
\CtwoCoverage{\mathrm{Cov}}(\CtwoConfidence{90\%})&=\frac{\CtwoCoverage{133}}{\CtwoCoverage{150}}\times100\%\approx\CtwoCoverage{88.67\%};\\
\CtwoError{E_{\rm abs}}(\CtwoConfidence{90\%})&=\left|\frac{\CtwoCoverage{13300}}{\CtwoCoverage{150}}-\CtwoConfidence{90}\right|\approx\CtwoError{1.33}\ \text{percentage points}.
\end{align*}
Here $\CtwoCoverage{133}$ of $\CtwoCoverage{150}$ subsequent intervals cover the district mean monthly pet-care spending for all registered households in the specified month: observed coverage is below nominal \CtwoConfidence{90\%} confidence by \CtwoError{1.33} \CtwoError{percentage points}.
\Needspace{10\baselineskip}
\textbf{$\CtwoConfidence{99\%}$ confidence}
\begin{align*}
\CtwoCoverCount(\CtwoConfidence{99\%})&=\CtwoCoverage{150}-\CtwoCoverage{3}=\CtwoCoverage{147},\qquad \CtwoTotalCount(\CtwoConfidence{99\%})=\CtwoCoverage{150};\\
\CtwoCoverage{\mathrm{Cov}}(\CtwoConfidence{99\%})&=\frac{\CtwoCoverage{147}}{\CtwoCoverage{150}}\times100\%=\CtwoCoverage{98\%};\\
\CtwoError{E_{\rm abs}}(\CtwoConfidence{99\%})&=\left|\CtwoCoverage{98}-\CtwoConfidence{99}\right|=\CtwoError{1}\ \text{percentage points}.
\end{align*}
Here $\CtwoCoverage{147}$ of $\CtwoCoverage{150}$ subsequent intervals cover the district mean monthly pet-care spending for all registered households in the specified month: observed coverage is below nominal \CtwoConfidence{99\%} confidence by \CtwoError{1} \CtwoError{percentage points}.
For each study the intervals share its sample mean and \CtwoSE{SE}. Increasing \CtwoConfidence{confidence} increases $\CtwoCritical{z^*}$, \CtwoMoE{MoE}, and \CtwoWidth{full width}. The narrower interval lies inside each wider partner, so covering counts cannot decrease with confidence on these same studies. 
These are retrospective studies of the same fixed population and target period. Finite-run \CtwoCoverage{observed coverage} need not equal \CtwoConfidence{nominal confidence} or approach it monotonically as studies are added. Under the model, confidence describes long-run coverage of the sampling procedure, including the one used for the initial report. Each realized interval either covers or misses the fixed mean; coverage is not the fraction of individual observations inside the initial interval. A smaller finite-run \CtwoError{absolute error} does not establish superior long-run performance.
\clearpage
\subsection*{5. Studies of clinic waiting times}
\subsubsection*{Task 1 solution. Organize the information and calculate SE}
\SolutionInformationTable{
$\CtwoMean{\mu}$ & \CtwoMean{Population mean} & Population parameters & The mean waiting time for all clinic appointments in the specified month, in minutes; fixed and initially unknown. \\[6pt]
$\CtwoMean{\mu_{\mathrm{known}}}$ & \CtwoMean{Later established} \CtwoMean{population mean} &  & $\CtwoMean{27}$ minutes; established by the later comprehensive population investigation. \\[6pt]
$\CtwoSD{\sigma}$ & \CtwoSD{Population} \CtwoSD{standard deviation} &  & $\CtwoSD{14}$ minutes; known spread of individual observations. \\[6pt]
\midrule
$\CtwoPoint{\bar{x}}$ & \CtwoPoint{Sample mean} & Sample statistics & $\CtwoPoint{26}$ minutes; the initial point estimate of $\CtwoMean{\mu}$. \\[6pt]
$\CtwoSize{n}$ & \CtwoSize{Sample size} &  & $\CtwoSize{49}$ appointments in the initial sample and each subsequent study. \\[6pt]
$\CtwoSE{\mathrm{SE}}$ & \CtwoSE{Standard error} &  & $\CtwoSE{\mathrm{SE}}=\CtwoSD{\sigma}/\sqrt{\CtwoSize{n}}$\par$=\CtwoSD{14}/\sqrt{\CtwoSize{49}}=\CtwoSE{2}$ minutes. Spread of means across repeated independent samples of 49 appointments. \\[6pt]
\midrule
$\CtwoConfidence{C}$ & \CtwoConfidence{Confidence} \CtwoConfidence{level} & Procedure Settings & $\CtwoConfidence{90\%}$: $\CtwoConfidence{C}=\CtwoConfidence{0.90}$. \\[6pt]
$\CtwoSignificance{\alpha}$ & \CtwoSignificance{Significance} \CtwoSignificance{level} &  & $\CtwoConfidence{90\%}$: $\CtwoSignificance{\alpha}=1-\CtwoConfidence{0.90}=\CtwoSignificance{0.10}$. \\[6pt]
$\CtwoSignificance{\alpha/2}$ & \CtwoSignificance{Tail} \CtwoSignificance{area} &  & $\CtwoConfidence{90\%}$: $\CtwoSignificance{\alpha/2}=\CtwoSignificance{0.10}/2=\CtwoSignificance{0.05}$. \\[6pt]
$\CtwoCritical{z^*}$ & \CtwoCritical{Critical} \CtwoCritical{value} &  & $\CtwoConfidence{90\%}$: $\CtwoCritical{z^*}=\CtwoCritical{1.645}$. \\[6pt]
}
The observations are independent and random, $\CtwoSD{\sigma}$ is known, and the sample mean is normally distributed under the stated model. No sample SD is required.
\clearpage
\subsubsection*{Task 2 solution. Calculate MoE and state the interval}
\textbf{$\CtwoConfidence{90\%}$ confidence}
\SolutionTaskTwoTable{
$\CtwoPoint{\bar{x}}$ & \CtwoPoint{Point estimate} & $\CtwoPoint{\bar{x}}=\sum x_i/\CtwoSize{n}$\par Given sample mean & $\CtwoPoint{26}$\par minutes & Estimates the population mean. \\[6pt]
\midrule
$\CtwoMoE{\mathrm{MoE}}$ & \CtwoMoE{Margin of error} & $\CtwoMoE{\mathrm{MoE}}=\CtwoCritical{z^*}\times\CtwoSE{\mathrm{SE}}$\par$=\CtwoCritical{1.645}\times\CtwoSE{2}$ & $\CtwoMoE{3.29}$\par minutes & Distance from centre to either bound. \\[6pt]
$\CtwoEndpoint{L}$ & \CtwoEndpoint{Lower} \CtwoEndpoint{endpoint} & $\CtwoEndpoint{L}=\CtwoPoint{\bar{x}}-\CtwoMoE{\mathrm{MoE}}$\par$=\CtwoPoint{26}-\CtwoMoE{3.29}$ & $\CtwoEndpoint{22.71}$\par minutes & Lower bound for the population mean. \\[6pt]
$\CtwoEndpoint{U}$ & \CtwoEndpoint{Upper} \CtwoEndpoint{endpoint} & $\CtwoEndpoint{U}=\CtwoPoint{\bar{x}}+\CtwoMoE{\mathrm{MoE}}$\par$=\CtwoPoint{26}+\CtwoMoE{3.29}$ & $\CtwoEndpoint{29.29}$\par minutes & Upper bound for the population mean. \\[6pt]
$\CtwoCI{\mathrm{CI}}$ & \CtwoCI{Confidence} \CtwoCI{interval} & $\CtwoCI{\mathrm{CI}}=[\CtwoEndpoint{L},\CtwoEndpoint{U}]$\par$=[\CtwoPoint{26}-\CtwoMoE{3.29},$\par$\CtwoPoint{26}+\CtwoMoE{3.29}]$ & $[\CtwoEndpoint{22.71},$\par$\CtwoEndpoint{29.29}]$\par minutes & Plausible values for the population mean under the model. \\[6pt]
$\CtwoWidth{W}$ & \CtwoWidth{Full width} & $\CtwoWidth{W}=\CtwoEndpoint{U}-\CtwoEndpoint{L}$\par$=\CtwoEndpoint{29.29}-\CtwoEndpoint{22.71}$\par$=2\CtwoMoE{\mathrm{MoE}}=2(\CtwoMoE{3.29})$ & $\CtwoWidth{6.58}$\par minutes & Total span; smaller width means greater precision. \\[6pt]
}
We are \CtwoConfidence{90\%} confident that the mean waiting time for all clinic appointments in the specified month lies between $\CtwoEndpoint{22.71}$ and $\CtwoEndpoint{29.29}$ minutes. Under the model, \CtwoConfidence{90\%} of intervals constructed by this procedure would cover the fixed $\CtwoMean{\mu}$ in the long run. This is not a probability assigned to the fixed mean after calculating this interval.
\clearpage
\subsubsection*{Task 3 solution. Calculate coverage, absolute error, and interpret confidence}
For these subsequent studies, an interval covers exactly when $\CtwoEndpoint{L_i}\leq\CtwoMean{\mu_{\mathrm{known}}}\leq\CtwoEndpoint{U_i}$, including equality. The later established mean is $\CtwoMean{27}$ minutes. Each interval is centred on its own study sample mean.
The counts by study-order block are $\CtwoCoverage{4,5,4,3}$. Add them cumulatively to assess how observed coverage changes as studies are added.
\Needspace{10\baselineskip}
\textbf{$\CtwoConfidence{90\%}$ confidence -- first 5 studies}
\begin{align*}
\CtwoCoverCount(\CtwoConfidence{90\%})&=\CtwoCoverage{4}=\CtwoCoverage{4},\qquad \CtwoTotalCount(\CtwoConfidence{90\%})=\CtwoCoverage{5};\\
\CtwoCoverage{\mathrm{Cov}}(\CtwoConfidence{90\%})&=\frac{\CtwoCoverage{4}}{\CtwoCoverage{5}}\times100\%=\CtwoCoverage{80\%};\\
\CtwoError{E_{\rm abs}}(\CtwoConfidence{90\%})&=\left|\CtwoCoverage{80}-\CtwoConfidence{90}\right|=\CtwoError{10}\ \text{percentage points}.
\end{align*}
Here $\CtwoCoverage{4}$ of $\CtwoCoverage{5}$ subsequent intervals cover the mean waiting time for all clinic appointments in the specified month: observed coverage is below nominal \CtwoConfidence{90\%} confidence by \CtwoError{10} \CtwoError{percentage points}.
\Needspace{10\baselineskip}
\textbf{$\CtwoConfidence{90\%}$ confidence -- first 10 studies}
\begin{align*}
\CtwoCoverCount(\CtwoConfidence{90\%})&=\CtwoCoverage{4}+\CtwoCoverage{5}=\CtwoCoverage{9},\qquad \CtwoTotalCount(\CtwoConfidence{90\%})=\CtwoCoverage{10};\\
\CtwoCoverage{\mathrm{Cov}}(\CtwoConfidence{90\%})&=\frac{\CtwoCoverage{9}}{\CtwoCoverage{10}}\times100\%=\CtwoCoverage{90\%};\\
\CtwoError{E_{\rm abs}}(\CtwoConfidence{90\%})&=\left|\CtwoCoverage{90}-\CtwoConfidence{90}\right|=\CtwoError{0}\ \text{percentage points}.
\end{align*}
Here $\CtwoCoverage{9}$ of $\CtwoCoverage{10}$ subsequent intervals cover the mean waiting time for all clinic appointments in the specified month: observed coverage is equal to nominal \CtwoConfidence{90\%} confidence.
\Needspace{10\baselineskip}
\textbf{$\CtwoConfidence{90\%}$ confidence -- first 15 studies}
\begin{align*}
\CtwoCoverCount(\CtwoConfidence{90\%})&=\CtwoCoverage{4}+\CtwoCoverage{5}+\CtwoCoverage{4}=\CtwoCoverage{13},\qquad \CtwoTotalCount(\CtwoConfidence{90\%})=\CtwoCoverage{15};\\
\CtwoCoverage{\mathrm{Cov}}(\CtwoConfidence{90\%})&=\frac{\CtwoCoverage{13}}{\CtwoCoverage{15}}\times100\%\approx\CtwoCoverage{86.67\%};\\
\CtwoError{E_{\rm abs}}(\CtwoConfidence{90\%})&=\left|\frac{\CtwoCoverage{1300}}{\CtwoCoverage{15}}-\CtwoConfidence{90}\right|\approx\CtwoError{3.33}\ \text{percentage points}.
\end{align*}
Here $\CtwoCoverage{13}$ of $\CtwoCoverage{15}$ subsequent intervals cover the mean waiting time for all clinic appointments in the specified month: observed coverage is below nominal \CtwoConfidence{90\%} confidence by \CtwoError{3.33} \CtwoError{percentage points}.
\clearpage
\subsubsection*{Task 3 solution (continued). Final coverage and interpretation}
\Needspace{10\baselineskip}
\textbf{$\CtwoConfidence{90\%}$ confidence -- first 20 studies}
\begin{align*}
\CtwoCoverCount(\CtwoConfidence{90\%})&=\CtwoCoverage{4}+\CtwoCoverage{5}+\CtwoCoverage{4}+\CtwoCoverage{3}=\CtwoCoverage{16},\qquad \CtwoTotalCount(\CtwoConfidence{90\%})=\CtwoCoverage{20};\\
\CtwoCoverage{\mathrm{Cov}}(\CtwoConfidence{90\%})&=\frac{\CtwoCoverage{16}}{\CtwoCoverage{20}}\times100\%=\CtwoCoverage{80\%};\\
\CtwoError{E_{\rm abs}}(\CtwoConfidence{90\%})&=\left|\CtwoCoverage{80}-\CtwoConfidence{90}\right|=\CtwoError{10}\ \text{percentage points}.
\end{align*}
Here $\CtwoCoverage{16}$ of $\CtwoCoverage{20}$ subsequent intervals cover the mean waiting time for all clinic appointments in the specified month: observed coverage is below nominal \CtwoConfidence{90\%} confidence by \CtwoError{10} \CtwoError{percentage points}.
The cumulative rates are $\CtwoCoverage{80\%,90\%,86.67\%,80\%}$ (the third is rounded); absolute errors are $\CtwoError{10,0,3.33,10}$ \CtwoError{percentage points}. Thus the rate first reaches nominal confidence and then moves away. The final covering count is $\CtwoCoverage{4}+\CtwoCoverage{5}+\CtwoCoverage{4}+\CtwoCoverage{3}=\CtwoCoverage{16}$ out of $\CtwoCoverage{20}$. More studies do not guarantee monotonic convergence.
These are retrospective studies of the same fixed population and target period. Finite-run \CtwoCoverage{observed coverage} need not equal \CtwoConfidence{nominal confidence} or approach it monotonically as studies are added. Under the model, confidence describes long-run coverage of the sampling procedure, including the one used for the initial report. Each realized interval either covers or misses the fixed mean; coverage is not the fraction of individual observations inside the initial interval. A smaller finite-run \CtwoError{absolute error} does not establish superior long-run performance.
\clearpage
\subsection*{6. Three confidence levels for library processing times}
\subsubsection*{Task 1 solution. Organize the information and calculate SE}
\SolutionInformationTable{
$\CtwoMean{\mu}$ & \CtwoMean{Population mean} & Population parameters & The mean processing time for all books returned to the library in the specified month, in minutes; fixed and initially unknown. \\[6pt]
$\CtwoMean{\mu_{\mathrm{known}}}$ & \CtwoMean{Later established} \CtwoMean{population mean} &  & $\CtwoMean{25}$ minutes; established by the later comprehensive population investigation. \\[6pt]
$\CtwoSD{\sigma}$ & \CtwoSD{Population} \CtwoSD{standard deviation} &  & $\CtwoSD{15}$ minutes; known spread of individual observations. \\[6pt]
\midrule
$\CtwoPoint{\bar{x}}$ & \CtwoPoint{Sample mean} & Sample statistics & $\CtwoPoint{24}$ minutes; the initial point estimate of $\CtwoMean{\mu}$. \\[6pt]
$\CtwoSize{n}$ & \CtwoSize{Sample size} &  & $\CtwoSize{100}$ books in the initial sample and each subsequent study. \\[6pt]
$\CtwoSE{\mathrm{SE}}$ & \CtwoSE{Standard error} &  & $\CtwoSE{\mathrm{SE}}=\CtwoSD{\sigma}/\sqrt{\CtwoSize{n}}$\par$=\CtwoSD{15}/\sqrt{\CtwoSize{100}}=\CtwoSE{1.5}$ minutes. Spread of means across repeated independent samples of 100 books. \\[6pt]
\midrule
$\CtwoConfidence{C}$ & \CtwoConfidence{Confidence} \CtwoConfidence{level} & Procedure Settings & $\CtwoConfidence{90\%}$: $\CtwoConfidence{C}=\CtwoConfidence{0.90}$\par $\CtwoConfidence{95\%}$: $\CtwoConfidence{C}=\CtwoConfidence{0.95}$\par $\CtwoConfidence{99\%}$: $\CtwoConfidence{C}=\CtwoConfidence{0.99}$. \\[6pt]
$\CtwoSignificance{\alpha}$ & \CtwoSignificance{Significance} \CtwoSignificance{level} &  & $\CtwoConfidence{90\%}$: $\CtwoSignificance{\alpha}=1-\CtwoConfidence{0.90}=\CtwoSignificance{0.10}$\par $\CtwoConfidence{95\%}$: $\CtwoSignificance{\alpha}=1-\CtwoConfidence{0.95}=\CtwoSignificance{0.05}$\par $\CtwoConfidence{99\%}$: $\CtwoSignificance{\alpha}=1-\CtwoConfidence{0.99}=\CtwoSignificance{0.01}$. \\[6pt]
$\CtwoSignificance{\alpha/2}$ & \CtwoSignificance{Tail} \CtwoSignificance{area} &  & $\CtwoConfidence{90\%}$: $\CtwoSignificance{\alpha/2}=\CtwoSignificance{0.10}/2=\CtwoSignificance{0.05}$\par $\CtwoConfidence{95\%}$: $\CtwoSignificance{\alpha/2}=\CtwoSignificance{0.05}/2=\CtwoSignificance{0.025}$\par $\CtwoConfidence{99\%}$: $\CtwoSignificance{\alpha/2}=\CtwoSignificance{0.01}/2=\CtwoSignificance{0.005}$. \\[6pt]
$\CtwoCritical{z^*}$ & \CtwoCritical{Critical} \CtwoCritical{value} &  & $\CtwoConfidence{90\%}$: $\CtwoCritical{z^*}=\CtwoCritical{1.645}$\par $\CtwoConfidence{95\%}$: $\CtwoCritical{z^*}=\CtwoCritical{1.960}$\par $\CtwoConfidence{99\%}$: $\CtwoCritical{z^*}=\CtwoCritical{2.576}$. \\[6pt]
}
The observations are independent and random, $\CtwoSD{\sigma}$ is known, and the sample mean is normally distributed under the stated model. No sample SD is required.
\clearpage
\subsubsection*{Task 2 solution. Calculate MoE and state the interval}
\textbf{$\CtwoConfidence{90\%}$ confidence}
\SolutionTaskTwoTable{
$\CtwoPoint{\bar{x}}$ & \CtwoPoint{Point estimate} & $\CtwoPoint{\bar{x}}=\sum x_i/\CtwoSize{n}$\par Given sample mean & $\CtwoPoint{24}$\par minutes & Estimates the population mean. \\[6pt]
\midrule
$\CtwoMoE{\mathrm{MoE}}$ & \CtwoMoE{Margin of error} & $\CtwoMoE{\mathrm{MoE}}=\CtwoCritical{z^*}\times\CtwoSE{\mathrm{SE}}$\par$=\CtwoCritical{1.645}\times\CtwoSE{1.5}$ & $\CtwoMoE{2.4675}$\par minutes & Distance from centre to either bound. \\[6pt]
$\CtwoEndpoint{L}$ & \CtwoEndpoint{Lower} \CtwoEndpoint{endpoint} & $\CtwoEndpoint{L}=\CtwoPoint{\bar{x}}-\CtwoMoE{\mathrm{MoE}}$\par$=\CtwoPoint{24}-\CtwoMoE{2.4675}$ & $\CtwoEndpoint{21.5325}$\par minutes & Lower bound for the population mean. \\[6pt]
$\CtwoEndpoint{U}$ & \CtwoEndpoint{Upper} \CtwoEndpoint{endpoint} & $\CtwoEndpoint{U}=\CtwoPoint{\bar{x}}+\CtwoMoE{\mathrm{MoE}}$\par$=\CtwoPoint{24}+\CtwoMoE{2.4675}$ & $\CtwoEndpoint{26.4675}$\par minutes & Upper bound for the population mean. \\[6pt]
$\CtwoCI{\mathrm{CI}}$ & \CtwoCI{Confidence} \CtwoCI{interval} & $\CtwoCI{\mathrm{CI}}=[\CtwoEndpoint{L},\CtwoEndpoint{U}]$\par$=[\CtwoPoint{24}-\CtwoMoE{2.4675},$\par$\CtwoPoint{24}+\CtwoMoE{2.4675}]$ & $[\CtwoEndpoint{21.5325},$\par$\CtwoEndpoint{26.4675}]$\par minutes & Plausible values for the population mean under the model. \\[6pt]
$\CtwoWidth{W}$ & \CtwoWidth{Full width} & $\CtwoWidth{W}=\CtwoEndpoint{U}-\CtwoEndpoint{L}$\par$=\CtwoEndpoint{26.4675}-\CtwoEndpoint{21.5325}$\par$=2\CtwoMoE{\mathrm{MoE}}=2(\CtwoMoE{2.4675})$ & $\CtwoWidth{4.935}$\par minutes & Total span; smaller width means greater precision. \\[6pt]
}
We are \CtwoConfidence{90\%} confident that the mean processing time for all books returned to the library in the specified month lies between $\CtwoEndpoint{21.5325}$ and $\CtwoEndpoint{26.4675}$ minutes. Under the model, \CtwoConfidence{90\%} of intervals constructed by this procedure would cover the fixed $\CtwoMean{\mu}$ in the long run. This is not a probability assigned to the fixed mean after calculating this interval.
\clearpage
\subsubsection*{Task 2 solution (continued). Calculate MoE and state the interval}
\textbf{$\CtwoConfidence{95\%}$ confidence}
\SolutionTaskTwoTable{
\multicolumn{5}{@{}l}{Shared \CtwoPoint{point estimate}: $\CtwoPoint{\bar{x}}$=$\CtwoPoint{24}$ minutes.}\\[6pt]
\midrule
$\CtwoMoE{\mathrm{MoE}}$ & \CtwoMoE{Margin of error} & $\CtwoMoE{\mathrm{MoE}}=\CtwoCritical{z^*}\times\CtwoSE{\mathrm{SE}}$\par$=\CtwoCritical{1.960}\times\CtwoSE{1.5}$ & $\CtwoMoE{2.94}$\par minutes & Distance from centre to either bound. \\[6pt]
$\CtwoEndpoint{L}$ & \CtwoEndpoint{Lower} \CtwoEndpoint{endpoint} & $\CtwoEndpoint{L}=\CtwoPoint{\bar{x}}-\CtwoMoE{\mathrm{MoE}}$\par$=\CtwoPoint{24}-\CtwoMoE{2.94}$ & $\CtwoEndpoint{21.06}$\par minutes & Lower bound for the population mean. \\[6pt]
$\CtwoEndpoint{U}$ & \CtwoEndpoint{Upper} \CtwoEndpoint{endpoint} & $\CtwoEndpoint{U}=\CtwoPoint{\bar{x}}+\CtwoMoE{\mathrm{MoE}}$\par$=\CtwoPoint{24}+\CtwoMoE{2.94}$ & $\CtwoEndpoint{26.94}$\par minutes & Upper bound for the population mean. \\[6pt]
$\CtwoCI{\mathrm{CI}}$ & \CtwoCI{Confidence} \CtwoCI{interval} & $\CtwoCI{\mathrm{CI}}=[\CtwoEndpoint{L},\CtwoEndpoint{U}]$\par$=[\CtwoPoint{24}-\CtwoMoE{2.94},$\par$\CtwoPoint{24}+\CtwoMoE{2.94}]$ & $[\CtwoEndpoint{21.06},$\par$\CtwoEndpoint{26.94}]$\par minutes & Plausible values for the population mean under the model. \\[6pt]
$\CtwoWidth{W}$ & \CtwoWidth{Full width} & $\CtwoWidth{W}=\CtwoEndpoint{U}-\CtwoEndpoint{L}$\par$=\CtwoEndpoint{26.94}-\CtwoEndpoint{21.06}$\par$=2\CtwoMoE{\mathrm{MoE}}=2(\CtwoMoE{2.94})$ & $\CtwoWidth{5.88}$\par minutes & Total span; smaller width means greater precision. \\[6pt]
}
We are \CtwoConfidence{95\%} confident that the mean processing time for all books returned to the library in the specified month lies between $\CtwoEndpoint{21.06}$ and $\CtwoEndpoint{26.94}$ minutes. Under the model, \CtwoConfidence{95\%} of intervals constructed by this procedure would cover the fixed $\CtwoMean{\mu}$ in the long run. This is not a probability assigned to the fixed mean after calculating this interval.
\clearpage
\subsubsection*{Task 2 solution (continued). Calculate MoE and state the interval}
\textbf{$\CtwoConfidence{99\%}$ confidence}
\SolutionTaskTwoTable{
\multicolumn{5}{@{}l}{Shared \CtwoPoint{point estimate}: $\CtwoPoint{\bar{x}}$=$\CtwoPoint{24}$ minutes.}\\[6pt]
\midrule
$\CtwoMoE{\mathrm{MoE}}$ & \CtwoMoE{Margin of error} & $\CtwoMoE{\mathrm{MoE}}=\CtwoCritical{z^*}\times\CtwoSE{\mathrm{SE}}$\par$=\CtwoCritical{2.576}\times\CtwoSE{1.5}$ & $\CtwoMoE{3.864}$\par minutes & Distance from centre to either bound. \\[6pt]
$\CtwoEndpoint{L}$ & \CtwoEndpoint{Lower} \CtwoEndpoint{endpoint} & $\CtwoEndpoint{L}=\CtwoPoint{\bar{x}}-\CtwoMoE{\mathrm{MoE}}$\par$=\CtwoPoint{24}-\CtwoMoE{3.864}$ & $\CtwoEndpoint{20.136}$\par minutes & Lower bound for the population mean. \\[6pt]
$\CtwoEndpoint{U}$ & \CtwoEndpoint{Upper} \CtwoEndpoint{endpoint} & $\CtwoEndpoint{U}=\CtwoPoint{\bar{x}}+\CtwoMoE{\mathrm{MoE}}$\par$=\CtwoPoint{24}+\CtwoMoE{3.864}$ & $\CtwoEndpoint{27.864}$\par minutes & Upper bound for the population mean. \\[6pt]
$\CtwoCI{\mathrm{CI}}$ & \CtwoCI{Confidence} \CtwoCI{interval} & $\CtwoCI{\mathrm{CI}}=[\CtwoEndpoint{L},\CtwoEndpoint{U}]$\par$=[\CtwoPoint{24}-\CtwoMoE{3.864},$\par$\CtwoPoint{24}+\CtwoMoE{3.864}]$ & $[\CtwoEndpoint{20.136},$\par$\CtwoEndpoint{27.864}]$\par minutes & Plausible values for the population mean under the model. \\[6pt]
$\CtwoWidth{W}$ & \CtwoWidth{Full width} & $\CtwoWidth{W}=\CtwoEndpoint{U}-\CtwoEndpoint{L}$\par$=\CtwoEndpoint{27.864}-\CtwoEndpoint{20.136}$\par$=2\CtwoMoE{\mathrm{MoE}}=2(\CtwoMoE{3.864})$ & $\CtwoWidth{7.728}$\par minutes & Total span; smaller width means greater precision. \\[6pt]
}
We are \CtwoConfidence{99\%} confident that the mean processing time for all books returned to the library in the specified month lies between $\CtwoEndpoint{20.136}$ and $\CtwoEndpoint{27.864}$ minutes. Under the model, \CtwoConfidence{99\%} of intervals constructed by this procedure would cover the fixed $\CtwoMean{\mu}$ in the long run. This is not a probability assigned to the fixed mean after calculating this interval.
\clearpage
\subsubsection*{Task 3 solution. Calculate coverage, absolute error, and interpret confidence}
For these subsequent studies, an interval covers exactly when $\CtwoEndpoint{L_i}\leq\CtwoMean{\mu_{\mathrm{known}}}\leq\CtwoEndpoint{U_i}$, including equality. The later established mean is $\CtwoMean{25}$ minutes. Each interval is centred on its own study sample mean.
The mutually exclusive categories total $\CtwoCoverage{219}+\CtwoCoverage{18}+\CtwoCoverage{10}+\CtwoCoverage{3}=\CtwoCoverage{250}$. Sum the relevant categories for each level; the denominator is $\CtwoCoverage{250}$ studies at each level.
\Needspace{10\baselineskip}
\textbf{$\CtwoConfidence{90\%}$ confidence}
\begin{align*}
\CtwoCoverCount(\CtwoConfidence{90\%})&=\CtwoCoverage{219}=\CtwoCoverage{219},\qquad \CtwoTotalCount(\CtwoConfidence{90\%})=\CtwoCoverage{250};\\
\CtwoCoverage{\mathrm{Cov}}(\CtwoConfidence{90\%})&=\frac{\CtwoCoverage{219}}{\CtwoCoverage{250}}\times100\%\approx\CtwoCoverage{87.60\%};\\
\CtwoError{E_{\rm abs}}(\CtwoConfidence{90\%})&=\left|\frac{\CtwoCoverage{21900}}{\CtwoCoverage{250}}-\CtwoConfidence{90}\right|\approx\CtwoError{2.40}\ \text{percentage points}.
\end{align*}
Here $\CtwoCoverage{219}$ of $\CtwoCoverage{250}$ subsequent intervals cover the mean processing time for all books returned to the library in the specified month: observed coverage is below nominal \CtwoConfidence{90\%} confidence by \CtwoError{2.40} \CtwoError{percentage points}.
\Needspace{10\baselineskip}
\textbf{$\CtwoConfidence{95\%}$ confidence}
\begin{align*}
\CtwoCoverCount(\CtwoConfidence{95\%})&=\CtwoCoverage{219}+\CtwoCoverage{18}=\CtwoCoverage{237},\qquad \CtwoTotalCount(\CtwoConfidence{95\%})=\CtwoCoverage{250};\\
\CtwoCoverage{\mathrm{Cov}}(\CtwoConfidence{95\%})&=\frac{\CtwoCoverage{237}}{\CtwoCoverage{250}}\times100\%\approx\CtwoCoverage{94.80\%};\\
\CtwoError{E_{\rm abs}}(\CtwoConfidence{95\%})&=\left|\frac{\CtwoCoverage{23700}}{\CtwoCoverage{250}}-\CtwoConfidence{95}\right|\approx\CtwoError{0.20}\ \text{percentage points}.
\end{align*}
Here $\CtwoCoverage{237}$ of $\CtwoCoverage{250}$ subsequent intervals cover the mean processing time for all books returned to the library in the specified month: observed coverage is below nominal \CtwoConfidence{95\%} confidence by \CtwoError{0.20} \CtwoError{percentage points}.
\Needspace{10\baselineskip}
\textbf{$\CtwoConfidence{99\%}$ confidence}
\begin{align*}
\CtwoCoverCount(\CtwoConfidence{99\%})&=\CtwoCoverage{219}+\CtwoCoverage{18}+\CtwoCoverage{10}=\CtwoCoverage{247},\qquad \CtwoTotalCount(\CtwoConfidence{99\%})=\CtwoCoverage{250};\\
\CtwoCoverage{\mathrm{Cov}}(\CtwoConfidence{99\%})&=\frac{\CtwoCoverage{247}}{\CtwoCoverage{250}}\times100\%\approx\CtwoCoverage{98.80\%};\\
\CtwoError{E_{\rm abs}}(\CtwoConfidence{99\%})&=\left|\frac{\CtwoCoverage{24700}}{\CtwoCoverage{250}}-\CtwoConfidence{99}\right|\approx\CtwoError{0.20}\ \text{percentage points}.
\end{align*}
Here $\CtwoCoverage{247}$ of $\CtwoCoverage{250}$ subsequent intervals cover the mean processing time for all books returned to the library in the specified month: observed coverage is below nominal \CtwoConfidence{99\%} confidence by \CtwoError{0.20} \CtwoError{percentage points}.
For each study the intervals share its sample mean and \CtwoSE{SE}. Increasing \CtwoConfidence{confidence} increases $\CtwoCritical{z^*}$, \CtwoMoE{MoE}, and \CtwoWidth{full width}. The narrower interval lies inside each wider partner, so covering counts cannot decrease with confidence on these same studies. The categories respect nesting: a study in ``all three counts at every level; ``only 95\% and 99\% counts at those two levels; ``only 99\% counts once at 99\%; ``none counts at no level. 
These are retrospective studies of the same fixed population and target period. Finite-run \CtwoCoverage{observed coverage} need not equal \CtwoConfidence{nominal confidence} or approach it monotonically as studies are added. Under the model, confidence describes long-run coverage of the sampling procedure, including the one used for the initial report. Each realized interval either covers or misses the fixed mean; coverage is not the fraction of individual observations inside the initial interval. A smaller finite-run \CtwoError{absolute error} does not establish superior long-run performance.
\par\clearpage
\endgroup
\restoregeometry
\end{document}
